\section{An exact zero-cost value of the collision bit}\label{sec:mark-value}
The previous result ties the detector's polynomial rank to a sharp
memory--performance law. A separate, elementary but exact decision
comparison shows that the additional bit also has strictly positive value
when no memory restriction is imposed. The comparison is not the
acceptance-cost mechanism of Section~\ref{sec:advantage}.

Fix the four-symbol detector of Section~\ref{sec:calibrated}. Consume two
cartridges, both without censoring ($g=1$, permitted by
Theorem~\ref{thm:boundary}). After the first record announce $a\in[0,1]$.
Let $Y$ indicate placement failure on the second cartridge, and receive
$1-(a-Y)^2$. The second record's other coordinates are irrelevant to this
score. Compare the best measurable forecaster with the first collision
sign available to the best measurable forecaster with only the first
placement and hit tags. The same two ambient attempts are consumed in both
experiments; no reward or resource saving is attached to acceptance.

Let $\mu_j=\int R^j\pi_0(dR)$ and define the posterior on a first hit by
$\pi_H(dR)=R\pi_0(dR)/\mu_1$. Put
\[
 v=\operatorname{Var}_{\pi_H}(R^2),\qquad
 \bar c=\tfrac12+\frac\epsilon\pi\mathbb E_{\pi_H}R^2.
\]
\begin{theorem}[Strict observable prediction gain, with an exact formula]\label{thm:zero-cost-gain}
For every full-support prior on $I$ and every fixed $T>0$, $\epsilon>0$
satisfying the apparatus bounds, the increase in optimal expected reward
from retaining the collision sign is exactly
\begin{equation}\label{eq:zero-cost-gain}
 \Delta_{\rm bit}
 =\frac{2T\mu_1}{a_0}\,
   \frac{(\epsilon v/a_0)^2}{\bar c(1-\bar c)}
 >0.
\end{equation}
At $\epsilon=0$ it is zero. Erasing the sign leaves the placement and hit
probabilities unchanged and removes exactly the cubic direction from the
finite attainable factor space. Thus the same structural feature is
responsible for the change of decision-state dimension, the change of
memory exponent, and the positive gain in this score problem.
\end{theorem}
\begin{proof}
The first hit has probability $2T\mu_1/a_0$. Off this event the two
information sets are identical. Conditional on it, the sign is a Bernoulli
observation with success probability
$c(R)=1/2+\epsilon R^2/\pi$. The next placement-failure probability given
$R$ is $h(R)=\pi R^2/a_0$. Independence of the two cartridges conditional
on $R$ makes the optimal forecast the conditional mean of $h(R)$.
For a binary observation $S$ with mean $\bar c$, direct Bayes calculation
gives
\[
 \mathbb E[h\mid S=1]-\mathbb E h
     =\frac{\operatorname{Cov}(h,c)}{\bar c},\qquad
 \mathbb E[h\mid S=0]-\mathbb E h
     =-\frac{\operatorname{Cov}(h,c)}{1-\bar c}.
\]
Consequently the increase in the variance of the forecast is
$\operatorname{Cov}(h,c)^2/\{\bar c(1-\bar c)\}$. The conditional-variance
identity for squared prediction loss identifies this with its Bayes-risk
improvement. Under $\pi_H$ the covariance is
$\epsilon v/a_0$. Multiplying by the first-hit probability proves
\eqref{eq:zero-cost-gain}. Full support on a nondegenerate positive interval
implies $v>0$, and the detector bound gives $0<\bar c<1$. The zero-amplitude
and rank-erasure assertions follow from \eqref{eq:four-cells}.
\end{proof}

This calculation is an explicit strictness instance of the comparison of
experiments and conditional-variance principles \cite{Blackwell}; those
principles are not claimed as new. The gain is over all measurable terminal
forecasters for the two stated information sets, not over a sampled menu.
Its constant can be small for a narrowly concentrated prior. Neither a
large numerical gain nor a universal gain independent of the prior is
asserted. Together with the sharp law of Section~\ref{sec:memory-risk}, it
provides a structural ablation which the retained acceptance-cost example
alone does not supply.
